% ==============================================================================
\section{Solving Approaches}
\label{sec:approaches}
% ==============================================================================

The MILP of \S\ref{sec:milp} is exact but its size grows quickly with the
number of aircraft, the per-aircraft chain length, and the blocking-arc
density: the partition variables $z^-,\,z^+,\,b^B,\,b^C$ contribute
$O(|\calR|^2 \cdot |\calA| \cdot (N_{\max} + 1))$ binaries.  We
complement the exact solver with two heuristic building blocks --- a
topology-aware GRASP that decides only the position assignment and a
fixed-assignment scheduler that solves the timing problem exactly given
that assignment --- plus a safe-pipeline composition that combines them
without any additional optimisation.

% ------------------------------------------------------------------------------
\subsection{Topology-Aware GRASP}
\label{sec:topo_heuristic}
% ------------------------------------------------------------------------------

The construction step inherits the structural intuition of an
aircraft-as-a-block heuristic: heavy aircraft assigned to topologically
central positions force their long occupancy windows to overlap with
aircraft serviced behind them, triggering manoeuvres or large delays.
We measure topological centrality through the \emph{blocking load} of a
position,
\[
  \ell(p) = |\{q \in \calP : p \text{ reaches } q \text{ in } G\}|,
\]
i.e.\ the number of positions that lie downstream of $p$ in the
blocking-arc graph, and bias the construction towards placing heavy
aircraft at positions with low (ideally zero) blocking load.

The candidate cost combines four terms,
\[
  c(r, p) = W^M\,\hat m + W^D\,\hat v^D_r + W^S\,\hat n_{rp}
            + W^T \cdot T_r \cdot \ell(p),
\]
where $\hat m, \hat v^D_r, \hat n_{rp}$ are produced by a fast
job-level rebuild (described below), $W^T$ is the topology-penalty
weight, and $T_r = \sum_{j \in \calJ_r} D_j$ is the aggregate
processing time of aircraft $r$.  The heuristic constructs a feasible
solution by biased-random greedy insertion, and diversification is
obtained by restarting from independent random seeds.
Algorithm~\ref{alg:topo} states the full procedure.

\begin{algorithm}[htbp]
\DontPrintSemicolon
\KwIn{Instance $I$, weights $(W^M, W^D, W^S, W^T)$,
      RCL parameter $\alpha$, time budget $T$, number of starts $n_s$.}
\KwOut{Best feasible job-level solution $z^*$ found.}
$z^* \leftarrow +\infty$; $T_s \leftarrow T / n_s$\;
\For{$k = 1, \ldots, n_s$}{
  $\mathcal{U} \leftarrow \calR$ sorted by $T_r$ descending\;
  \While{$\mathcal{U} \neq \emptyset$}{
    evaluate $c(r, p)$ for every $(r, p) \in \mathcal{U} \times \calP$
    using the Mode-A rebuild below\;
    pick $(r, p)$ from the geometrically biased RCL of weight
    parameter $\alpha$, commit it; remove $r$ from $\mathcal{U}$\;
  }
  $z \leftarrow$ Mode-A rebuild of the current assignment\;
  \If{$z < z^*$}{$z^* \leftarrow z$}
}
\Return $z^*$\;
\caption{Topology-aware GRASP for the job-level problem
         with multi-start ($n_s$ restarts).}
\label{alg:topo}
\end{algorithm}

\paragraph{Mode-A rebuild.}
The rebuild is the only piece that differs in spirit from an
aircraft-level GRASP.  Given a partial position assignment
$\pi : \calR' \to \calP$ over the aircraft committed so far, it
schedules the aircraft in increasing earliest-start order and, for
each, lays the job chain back-to-back at its position.  Before
committing the chain it adjusts the start so that the resulting
schedule is feasible under the three-mode access conditions of
\S\ref{subsec:access}:
\begin{itemize}[leftmargin=2em, itemsep=0pt]
  \item if any already-placed \emph{front} aircraft $r$ has a stay
        $[s_r, f_r]$ that overlaps the new aircraft's chain, push the
        new start to $f_r + \eta$ (clean Mode~A z+);
  \item if any already-placed \emph{rear} aircraft has an access
        instant $\tau$ that falls inside a job of the new chain, push
        the new chain start past that rear's full stay
        ($\mathrm{finish}_{\mathrm{rear}} + \eta$, again clean Mode~A);
  \item each shift may expose a new overlap with another aircraft, so
        both passes are wrapped in a convergence loop that only ever
        moves the chain forward and therefore terminates in a bounded
        number of iterations.
\end{itemize}
After the convergence loop, every access instant of every other
aircraft is either strictly before or strictly after the new
aircraft's stay (Mode~A everywhere).  No Mode-B gap is inserted and no
Mode-C interruption is triggered by the construction --- those two
modes would require retiming already-committed jobs (Mode~B) or
solving a fixpoint between extension lengths and access classifications
(Mode~C).  In a forward greedy both lead to non-monotone state updates
that we choose not to handle at this stage; Mode~B and Mode~C are
exploited downstream by the fixed-assignment scheduler.

\paragraph{Implication.}
The Mode-A rebuild trades two solution-quality components for a
substantially simpler construction.  First, $\hat n_{rp} = 0$ for every
GRASP candidate, so the movement term of the cost reduces to
$W^S \cdot 0$ during construction; movements only appear in the
post-construction objective when the rebuild later yields a schedule
in which a downstream rear was forced into Mode~A by delaying past a
front.  Second, the construction never benefits from a Mode-C
interruption that would shorten makespan or delay at the price of one
manoeuvre.  Both effects are recovered when the assignment is passed to
the fixed-assignment scheduler below, which solves the timing problem
to optimality and naturally identifies profitable Mode-B and Mode-C
configurations.

A local-search portfolio analogous to the aircraft-level
GRASP's---single-move, 2-opt swap, intra-position reorder,
perturb-and-rebuild---can be plugged on top of the Mode-A rebuild
because the operators only manipulate the position assignment (or the
within-position service order) and delegate the timing computation to
the rebuild.  We do not include such a portfolio in the present cut
to keep the heuristic's behaviour easy to reason about; the empirical
study in \S\ref{sec:experiments} therefore reports a pure
multi-start GRASP without local search.

% ------------------------------------------------------------------------------
\subsection{Fixed-Assignment Scheduler (FAS)}
\label{sec:fas}
% ------------------------------------------------------------------------------

Given a fixed position assignment $\pi : \calR \to \calP$, the problem
of \S\ref{sec:problem} reduces to a pure scheduling problem: find per-job
start times $\{s_j\}_{j \in \calJ}$ and interruption counters
$\{\kappa_j\}_{j \in \calJ}$ that minimise the objective
\eqref{eq:milpobj} subject to the constraints
\eqref{eq:assign}--\eqref{eq:makespan} of \S\ref{sec:milp}, with every
$x_{rp}$ replaced by the constant $\mathbf{1}[\pi(r) = p]$.

Substituting $\pi$ collapses the formulation along two axes:
\begin{itemize}[leftmargin=2em, itemsep=0pt]
  \item the $O(|\calR|\,|\calP|)$ assignment binaries
        $x_{rp}$ disappear, together with the assignment
        constraint~\eqref{eq:assign};
  \item the same-position ordering binaries $q_{rr'p}$ are only kept
        between aircraft that share a position under $\pi$, replacing
        the $O(|\calR|^2 |\calP|)$ enumeration of \eqref{eq:ordLB}--%
        \eqref{eq:ordTight} by an enumeration over actually-shared
        positions;
  \item the placement-compatibility variable $y_{rr'pp'}$ becomes the
        constant $\mathbf{1}[\pi(r){=}p \wedge \pi(r'){=}p']$, so the
        partition constraint~\eqref{eq:partition} only needs to be
        instantiated for the pairs $(r, r')$ that actually share an
        arc $(\pi(r), \pi(r')) \in \calA$.  All other partition
        variables are dropped.
\end{itemize}
The Mode-B gap, Mode-C extension and counter constraints
\eqref{eq:bb_lb}--\eqref{eq:kappa} carry over verbatim restricted to the
active arcs, as do the objective auxiliaries.
Algorithm~\ref{alg:fas} summarises the procedure.  The reduced model is
encoded directly in \texttt{gurobipy}, bypassing the Pyomo layer.

\begin{algorithm}[htbp]
\DontPrintSemicolon
\KwIn{Instance $I$, fixed assignment $\pi$, time budget $T$.}
\KwOut{Best feasible job-level schedule found within $T$.}
identify active arcs
   $\mathcal{A}^{\pi} = \{(r, r') : (\pi(r), \pi(r')) \in \calA\}$\;
build reduced MILP with per-job variables $s_j, f_j, \kappa_j$,
   partition variables only for pairs in $\mathcal{A}^{\pi}$ and both
   access sides, makespan, delay and movement auxiliaries\;
solve with Gurobi under time limit $T$\;
recover the aircraft-level summary
   $s_r \!=\! s_{j^r_1},\, f_r \!=\! f_{j^r_{N_r}}$
   and the schedule dict\;
\Return resulting schedule\;
\caption{Fixed-Assignment Scheduler (FAS) for the job-level problem.}
\label{alg:fas}
\end{algorithm}

In contrast to the GRASP, the FAS \emph{does} exploit Mode~B and
Mode~C whenever they improve the objective, since the partition
constraint forces every access to be classified into one of the four
admissible states and the solver is free to select the cheapest
combination across all blocking pairs.  Empirically the FAS therefore
absorbs a substantial part of the construction's Mode-A conservatism:
on the densely-blocked benchmark configurations the resulting movement
count is non-zero whenever an interruptible job in the front would
have allowed a strictly cheaper schedule than the corresponding Mode-A
delay.  The model size is dominated by the $4\,|\mathcal{A}^{\pi}|\,
(N_{\max}+1)$ partition binaries, which is one to two orders of
magnitude smaller than the full MILP of \S\ref{sec:milp} on the same
instance.

% ------------------------------------------------------------------------------
\subsection{Safe Pipeline}
\label{sec:safe_pipeline}
% ------------------------------------------------------------------------------

The two preceding methods can be run in sequence: the GRASP produces a
position assignment, and FAS then solves the timing problem on that
assignment.  Because the FAS receives a fixed $\pi$ it cannot recover
from a poor assignment; conversely the GRASP's Mode-A construction
leaves headroom that the FAS can convert into shorter makespan or fewer
movements.  Which of the two ends up with the lower objective is
instance-dependent.  The \emph{safe pipeline} returns the best
objective achieved by any of the component methods:
\[
  z^* = \min_{h \in \mathcal{H}} z_h,
  \qquad \mathcal{H} = \{\texttt{topology\_ms6\_job},\ \texttt{fas\_on\_topo\_job}\},
\]
and is therefore feasibility-preserving: as long as one component
returns a compliant solution, the pipeline returns the best of them.
Since each component already stores its solution in a shared cache,
the pipeline performs no additional optimisation; its cost is a
single comparison.

\begin{proposition}
The safe-pipeline objective satisfies $z^* \le z_h$ for every
$h \in \mathcal{H}$; in particular, $z^* \le z_{\texttt{topology\_ms6\_job}}$
and $z^* \le z_{\texttt{fas\_on\_topo\_job}}$, so the pipeline never
degrades either component's result.
\end{proposition}

A further consequence is robustness: when the GRASP fails to produce a
job-level-compliant schedule on a dense topology --- the construction
is Mode-A only and may push aircraft past the time budget, or expose
the rare residual conflict that the convergence loop cannot resolve in
forward order --- the FAS-on-topo branch still returns a compliant
solution as long as the assignment is feasible at all (the MILP either
proves infeasibility of the assignment or returns a valid schedule).
The pipeline therefore inherits the exactness of the MILP at the cost
of the GRASP runtime, without paying the build time of the full
job-level MILP of \S\ref{sec:milp}.
